\documentclass[11pt]{article}
\usepackage[margin=1in]{geometry}
\usepackage[T1]{fontenc}
\usepackage{graphicx}
\usepackage{hyperref}
\usepackage{enumitem}
\setlength{\parskip}{0.5em}
\setlength{\parindent}{0pt}

\title{test-time compute and reasoning in large language models}
\author{arXiv GraphRAG --- phase-6 deep research}
\date{\today}

\begin{document}
\maketitle

\begin{abstract}
An excellent, high-fidelity review that masterfully integrates 2025+ frontier research and critical theoretical tensions.
\end{abstract}

\section{Scaling Intelligence Beyond Training: A Review of Test-Time Compute and Reasoning in Large Language Models}

The advancement of large language models (LLMs) has historically centered on the scaling of pre-training parameters and datasets. However, as the industry encounters the diminishing returns of static model weights, the focus is shifting toward how models can utilize additional computational resources during the inference phase to solve complex, multi-step problems. This transition from "fast" pattern matching to "slow" deliberate reasoning is critical for moving LLMs from simple text generators to robust reasoning engines. Understanding how to effectively allocate compute at test-time is no longer just an optimization problem; it is a fundamental requirement for achieving higher-order cognitive tasks and approaching general intelligence.

This review traces the progression of this field, beginning with the foundations of LLMs and early reasoning paradigms that set the stage for modern architectures. We explore the subsequent development of reasoning capabilities through techniques such as Chain-of-Thought (CoT), the implementation of verifiers, and the emergence of self-correction mechanisms. The narrative then moves to the pivotal shift toward test-time scaling and the formalization of inference-time compute, supported by systematic surveys and the emerging scaling laws that define the relationship between additional computation and accuracy. Finally, we examine the state-of-the-art through the lens of advanced search, agentic workflows, and multi-agent synergy, before concluding with the frontiers of optimization, efficiency, and domain-specific scaling.

\subsection{Foundations of LLMs and Early Reasoning Paradigms}

The evolution of large language models (LLMs) is characterized by a progression from architectural standardization to the discovery of scaling laws and, ultimately, the development of specialized reasoning paradigms. The foundational Transformer architecture established the groundwork for modern sequence modeling, initially demonstrated through high performance in machine translation (English-to-German and English-to-French) and English constituency parsing using BLEU scores as a primary metric \texttt{arXiv:1706.03762}. This architectural stability enabled the identification of fundamental scaling laws, where cross-entropy loss follows a power-law relationship with model size, dataset size, and training compute \texttt{arXiv:2001.08361}.

As models scaled, research shifted toward understanding emergent capabilities. The transition to massive parameter counts—up to 175 billion—revealed that scaling significantly improves task-agnostic, few-shot performance across domains such as 3-digit arithmetic and word unscrambling \texttt{arXiv:2005.14165}. This scaling behavior was further quantified through multitask accuracy assessments, showing that larger models like GPT-3 outperformed random chance by nearly 20 percentage points across 57 diverse tasks, including law, morality, and computer science \texttt{arXiv:2009.03300}. However, these early scaling efforts faced notable bottlenecks: models often struggled with near-random accuracy on socially important subjects \texttt{arXiv:2009.03300}, and methodological issues arose from training on large web corpora \texttt{arXiv:2005.14165}.

The literature subsequently diverged into specialized reasoning methodologies to address these limitations. One trajectory focused on prompting techniques to elicit latent reasoning, such as the "Let's think step by step" zero-shot Chain of Thought (CoT) approach, which improved performance in arithmetic, symbolic, and logical reasoning \texttt{arXiv:2205.11916}. To enhance the reliability of these reasoning traces, self-consistency was introduced as a search strategy, sampling diverse reasoning paths and selecting the most consistent answer via marginalization to improve arithmetic and commonsense accuracy \texttt{arXiv:2203.11171}.

A second trajectory integrated LLMs with external computational tools. Program-aided language models (PAL) addressed mathematical and algorithmic reasoning by leveraging program execution \texttt{arXiv:2211.10435}, while other research focused on program synthesis, where performance scaled log-linearly with model size but remained limited by an inability to predict program outputs for specific inputs \texttt{arXiv:2108.07732}. Similarly, models trained on code demonstrated high success rates in Python synthesis through repeated sampling, though they struggled with complex docstrings involving long chains of operations \texttt{arXiv:2107.03374}. Finally, the ReAct paradigm moved beyond static reasoning by interleaving reasoning traces with task-specific actions, enabling interactive decision-making in environments like WebShop and fact verification in Fever \texttt{arXiv:2210.03629}.

\subsection{Developing Reasoning Capabilities: CoT, Verifiers, and Self-Correction}

While training-time compute optimization—such as scaling model size and training tokens equally (\texttt{arXiv:2203.15556})—is fundamental to achieving high-level performance like that seen in GPT-4 (\texttt{arXiv:2303.08774}), recent research shifts focus toward maximizing reasoning through test-time strategies.

Early improvements focused on prompting and sampling strategies to improve Chain-of-Thought (CoT) reasoning. Self-consistency (\texttt{arXiv:2203.11171}) enhances arithmetic and commonsense reasoning by sampling a diverse set of reasoning paths and selecting the most consistent answer via marginalization. This can be extended through structural decomposition; for instance, least-to-most prompting (\texttt{arXiv:2205.10625}) breaks complex problems into simpler subproblems solved in sequence. This approach demonstrates superior generalization compared to standard CoT, such as increasing SCAN length split accuracy from 16\% to at least 99\% using 14 exemplars. For even more non-trivial planning or search tasks, Tree of Thoughts (\texttt{arXiv:2305.10601}) moves beyond linear paths by employing exploration over coherent "thoughts" via lookahead and backtracking. This provides a significant leap in capability, achieving a 74\% success rate in the Game of 24 compared to only 4\% using CoT.

A parallel trajectory involves using verifiers to prune incorrect paths. While some methods utilize step-aware verifiers with weighted voting (\texttt{arXiv:2206.02336}) or process- and outcome-based feedback (\texttt{arXiv:2211.14275}), more recent approaches focus on step-by-step verification to improve accuracy on mathematical datasets like MATH (\texttt{arXiv:2305.20050}). Math-Shepherd (\texttt{arXiv:2312.08935}) further integrates verification with reinforcement learning via step-by-step PPO, boosting Mistral-7B's GSM8K accuracy from 77.9\% to 84.1\%. However, a notable weakness of such methods is their heavy reliance on manual annotation (\texttt{arXiv:2312.08935}).

Finally, reasoning can be enhanced through iterative self-improvement loops. STaR (\texttt{arXiv:2203.14465}) utilizes bootstrapping, where models are iteratively fine-tuned on self-generated rationales that lead to correct answers, enabling successively more complex reasoning. Similarly, Self-Refine (\texttt{arXiv:2303.17651}) employs a single LLM to act as generator, refiner, and feedback provider, yielding an average absolute performance improvement of \textasciitilde{}20\% across seven diverse tasks compared to one-step generation. Whereas prompting and verification focus on selecting the best path from a set of outputs, these iterative methods focus on refining the quality of the reasoning process itself.

\subsection{The Shift to Test-Time Scaling and Inference-Time Compute}

Recent research indicates a paradigm shift in large language model (LLM) optimization, moving from a primary focus on parameter scaling toward maximizing test-time compute to enhance reasoning capabilities. This transition is characterized by diverse mechanisms, ranging from repeated sampling and search-based strategies to reinforcement learning (RL) and recurrent latent processing.

A foundational approach to scaling inference involves increasing the number of candidate solutions through repeated sampling. \texttt{arXiv:2407.21787} demonstrates that in domains like coding and formal proofs, coverage scales across four orders of magnitude, often following an exponentiated power law. However, this method faces significant bottlenecks, as majority voting and reward models tend to plateau after several hundred samples, failing to fully exploit larger sample budgets [\texttt{arXiv:2407.21787}]. To address these inefficiencies, researchers have explored more sophisticated search strategies. \texttt{arXiv:2408.03314} proposes searching against dense, process-based verifier reward models and adaptively updating the model's distribution, achieving a compute-optimal strategy that is >4x more efficient than best-of-N [\texttt{arXiv:2408.03314}]. Similarly, \texttt{arXiv:2408.00724} analyzes various inference strategies—including greedy search, weighted voting, and tree search—finding that advanced algorithms allow smaller models to outperform significantly larger ones; for instance, Llemma-7B with tree search exceeds the performance of Llemma-34B on the MATH benchmark [\texttt{arXiv:2408.00724}].

The literature suggests that test-time scaling can fundamentally alter the relationship between model size and performance. \texttt{arXiv:2502.06703} provides a striking example, demonstrating that compute-optimal test-time scaling (TTS) can enable a 1B LLM to surpass a 405B LLM on MATH-500 tasks [\texttt{arXiv:2502.06703}]. This aligns with findings in \texttt{arXiv:2408.03314}, where test-time compute outperformed a 14x larger model in FLOPs-matched evaluations [\texttt{arXiv:2408.03314}]. While some methods focus on external search, others emphasize internalizing reasoning through training. \texttt{arXiv:2412.16720} and \texttt{arXiv:2501.12948} highlight the use of large-scale RL to incentivize chain-of-thought reasoning in mathematics and STEM fields [\texttt{arXiv:2412.16720}, \texttt{arXiv:2501.12948}].

Alternative scaling mechanisms also emerge in the form of "budget forcing" and latent depth. \texttt{arXiv:2501.19393} proposes "s1," a simple method that forces longer thinking processes (e.g., by appending "Wait") to extrapolate performance on competition math [\texttt{arXiv:2501.19393}]. In contrast, \texttt{arXiv:2502.05171} explores a recurrent depth approach, unrolling a recurrent block in latent space to improve performance up to a computation load equivalent to 50 billion parameters [\texttt{arXiv:2502.05171}]. Despite these advancements, the effectiveness of these strategies remains highly sensitive to prompt difficulty [\texttt{arXiv:2408.03314}], and the limitations of simple voting-based scaling necessitate the move toward the more complex, compute-optimal search and RL-driven architectures currently under investigation [\texttt{arXiv:2407.21787}, \texttt{arXiv:2501.12948}].

\subsection{Systematic Surveys and Scaling Laws of Inference Compute}

Recent literature has shifted from training-centric scaling to a rigorous investigation of test-time scaling (TTS), exploring how additional compute during inference can enhance model performance across mathematics, coding, and scientific reasoning [\texttt{arXiv:2503.24235}, \texttt{arXiv:2511.14772}]. While several surveys categorize TTS mechanisms—such as lengthening generated scratchpads, repeated model calls, or structural search strategies like Tree-of-Thought and Branch-Solve-Merge [\texttt{arXiv:2504.00294}, \texttt{arXiv:2511.14772}]—the literature reveals significant tension regarding the predictability and efficiency of these scaling laws.

A primary point of contention involves the monotonicity of scaling benefits. On one hand, some research suggests that optimal TTS performance scales monotonically with the allocated compute budget [\texttt{arXiv:2512.02008}]. This is complemented by perspectives that view TTS through the lens of subproblem organization, where compute is allocated via sequential, parallel, or tree-structured topological strategies to improve predictive accuracy [\texttt{arXiv:2511.14772}]. However, this optimistic view is challenged by evidence of "scaling plateaus." For instance, in mathematical and scientific reasoning domains (e.g., AIME, MATH-500, GPQA), research indicates that diminishing returns occur once a specific saturation point or threshold in the scaling budget is reached [\texttt{arXiv:2505.20522}].

The complexity of the task itself serves as a critical moderator of scaling effectiveness. While TTS can enhance performance in STEM reasoning, calendar planning, and spatial tasks, the advantages are not uniform; they tend to diminish as problem complexity increases [\texttt{arXiv:2504.00294}]. In fact, simply increasing token counts does not inherently guarantee higher accuracy in highly challenging regimes, suggesting that "more" compute is not a panacea for complex reasoning [\texttt{arXiv:2504.00294}].

Furthermore, the literature identifies a disconnect between algorithmic scaling and system-level efficiency. While most studies focus on accuracy as the primary evaluation metric [\texttt{arXiv:2504.00294}, \texttt{arXiv:2511.14772}], a systemic critique suggests that current methods are overly focused on the compute-optimal Pareto-frontier. This approach neglects practical system-optimal performance, such as the critical relationship between compute and real-world metrics like latency and cost-per-token [\texttt{arXiv:2509.19645}]. Consequently, while TTS offers pathways to improved reasoning through mechanisms like speculative decoding or parallel/sequential scaling [\texttt{arXiv:2509.19645}, \texttt{arXiv:2505.20522}], the field faces significant bottlenecks ranging from saturation points in reasoning capabilities to the economic and temporal costs of inference-time computation [\texttt{arXiv:2505.20522}, \texttt{arXiv:2509.19645}].

\subsection{Advanced Search, Agents, and Multi-Agent Synergy}

Recent advancements in test-time scaling have shifted from simple sampling to sophisticated search and agentic frameworks. A foundational approach involves AlphaZero-like tree-search, which utilizes a learned value function to guide decoding and training across tasks involving planning and decision-making, capable of handling trees with a depth of 64 [\texttt{arXiv:2309.17179}]. This tree-search paradigm is further refined in hierarchical contexts; for instance, \texttt{arXiv:2507.05557} integrates Monte Carlo Tree Search (MCTS) with hierarchical retrieval-augmented in-context learning—using coarse-level templates and fine-level solution steps—and a Process Reward Model (PRM) to achieve a 16\% relative improvement in mathematical reasoning tasks using LLaMA-3.1-8B.

While tree-search optimizes individual reasoning trajectories, agentic frameworks focus on managing complex, multi-stage tasks. AgentTTS [\texttt{arXiv:2508.00890}] optimizes compute-optimal scaling by dynamically allocating model selection and budgets to heterogeneous subtasks through iterative feedback-driven interactions. This contrasts with the broader scaling strategies for language agents proposed in \texttt{arXiv:2506.12928}, which emphasizes parallel sampling, sequential revision, and diversifying rollouts to improve agent performance. Similarly, \texttt{arXiv:2509.25420} introduces a dual-phase search that separates reasoning into distinct planning and execution phases, using reward-guided pruning to redistribute sampling effort and reduce redundant computation in math and code generation.

The literature also explores iterative self-correction through linguistic feedback and episodic memory, as seen in the Reflexion framework [\texttt{arXiv:2303.11366}], which utilizes verbal reinforcement learning for sequential decision-making. Building upon these concepts of memory and iteration, TMAS [\texttt{arXiv:2605.10344}] advances the field through multi-agent synergy. Unlike previous methods that suffer from weak coordination or noisy historical information, TMAS employs specialized agents and hierarchical memories—specifically an "experience bank" and a "guideline bank"—to balance exploration and exploitation. This approach demonstrates significant scaling effectiveness, notably increasing the Qwen3-4B-Thinking-2507 Pass@1 from 53.5\% to 73.5\% on IMO-AnswerBench.

However, the transition to multi-agent synergy introduces new trade-offs. While TMAS provides a novel hybrid reward RL scheme to incentivize strategy novelty, it relies heavily on a "cold-start" dataset distilled from a teacher model (DeepSeek-V3.2) and incurs high computational overhead due to the necessity of running multiple specialized agents and frequent verification passes. Thus, while agentic and multi-agent frameworks offer more structured ways to scale test-time compute, they face a tension between the depth of reasoning and the efficiency of the underlying computational allocation.

\subsection{Frontiers: Optimization, Efficiency, and Domain-Specific Scaling}

Recent advancements in test-time compute (TTC) emphasize a transition from brute-force sampling to optimized, domain-aware search strategies. While traditional model scaling is often prioritized, TTC can surpass it in both accuracy and energy efficiency for complex reasoning tasks [\texttt{arXiv:2505.14733}], provided the bottleneck of output sequence length is managed [\texttt{arXiv:2505.14733}].

Efficiency in TTC is increasingly achieved through dynamic allocation rather than static strategies. For instance, combining Best-of-N with beam search allows for more favorable accuracy, cost, and wall-clock latency trade-offs compared to static methods [\texttt{arXiv:2509.09864}]. This efficiency is further enhanced by Scaling-Law Guided (SLG) search, which utilizes tail-guided reward estimation to achieve expected rewards with a polynomially smaller compute budget than standard Best-of-N [\texttt{arXiv:2602.01485}]. Moving beyond full-trajectory sampling, research into thought-level self-certainty maximization demonstrates that selecting continuations based on the model's internal predictive certainty can match or exceed self-consistency under comparable token budgets [\texttt{arXiv:2603.07159}]. A significant strength of this approach is that optimizing only the first 1–5 reasoning steps can yield peak accuracy [\texttt{arXiv:2603.07159}]; however, a notable weakness is its reliance on predefined reasoning delimiters, which may limit its generalizability to models lacking structured Chain-of-Thought formatting [\texttt{arXiv:2603.07159}].

However, the effectiveness of TTC is heavily constrained by the reasoning domain and the nature of the required information. In knowledge-intensive tasks, increasing TTC does not consistently improve accuracy and may actually exacerbate hallucinations [\texttt{arXiv:2509.06861}]. This is attributed to an information-theoretic limit where the model cannot generate more truth-related information than what is already encoded in its parameters [\texttt{arXiv:2509.06861}]. Similarly, in vision-language reasoning, the benefits are inconsistent; while closed-source models show gains, open-source VLMs exhibit unpredictable behavior, and iterative self-refinement can sometimes degrade performance [\texttt{arXiv:2512.11109}].

Scaling behavior also exhibits diminishing returns across different dimensions. While increasing the sample number \$K\$ improves accuracy in mathematical, coding, and theory-of-mind tasks, performance eventually plateaus at large \$K\$, leaving certain hard questions unsolved [\texttt{arXiv:2510.02611}]. Interestingly, scaling along the temperature dimension can mitigate some of these limits, providing an additional 7.3 points of accuracy over single-temperature test-time scaling [\texttt{arXiv:2510.02611}].

\subsection{Tensions and Open Problems}

A central debate in the literature concerns the optimal mechanism for enhancing model performance: whether intelligence is primarily a product of massive pre-training or can be effectively recovered through expanded inference-time computation. While early paradigms focused on scaling parameters and few-shot prompting to achieve emergent capabilities (\texttt{arXiv:2005.14165}), recent shifts suggest that massive scale alone may be insufficient for complex reasoning. This tension is exemplified by the argument that small, highly optimized models can outperform massive counterparts when test-time compute is properly managed (\texttt{arXiv:2502.06703}), a view that challenges the "large monkey" approach of simply repeating massive inference passes (\texttt{arXiv:2407.21787}). Furthermore, the efficacy of these scaling methods remains contested; some suggest that temperature-based sampling is a critical driver for test-time scaling (\texttt{arXiv:2510.02611}), while others argue that the gains from such repetition may be fundamentally limited by the underlying model's reasoning structure (\texttt{arXiv:2407.21787}).

A second tension exists regarding the architectural implementation of reasoning: the choice between structured, deliberative search and iterative, self-correcting feedback loops. One school of thought advocates for explicit, tree-based search structures to facilitate deliberate problem-solving (\texttt{arXiv:2305.10601}), moving beyond the simpler majority-voting methods used in self-consistency (\texttt{arXiv:2203.11171}). Conversely, an alternative approach emphasizes the power of iterative refinement, where models use self-feedback to correct errors in a closed loop (\texttt{arXiv:2303.17651}). This debate is complicated by the role of reinforcement learning; while some argue that instruction tuning and human feedback are the primary drivers of helpfulness (\texttt{arXiv:2203.02155}), others suggest that true agentic reasoning requires verbal reinforcement learning and externalized reflection to achieve autonomous error correction (\texttt{arXiv:2303.11366}).

Finally, there is a growing disagreement over the efficiency and sustainability of scaling compute at inference time. As reasoning models increasingly rely on "scaling evaluation-time compute" to bridge the gap between base capabilities and expert-level performance (\texttt{arXiv:2503.19877}), the practical costs are becoming a bottleneck. While some frameworks attempt to maximize reasoning depth through multi-agent synergy (\texttt{arXiv:2605.10344}), others raise urgent concerns regarding the massive energy consumption and latency overheads inherent in these extended computation cycles; \texttt{arXiv:2509.09864}). This has led to a push for "reasoning on a budget," where the focus shifts from maximizing raw performance to developing adaptive, controllable compute strategies that balance accuracy with real-time constraints (\texttt{arXiv:2507.02076}).

\subsubsection{Outlook}

The frontier of the field is rapidly shifting from "scaling laws for training" to "scaling laws for inference." As we move toward 2026, the focus is converging on the discovery of optimal test-time scaling laws through reward-tail-guided searches (\texttt{arXiv:2602.01485}) and the minimization of uncertainty during the reasoning process (\texttt{arXiv:2603.07159}). The next generation of large language models will likely not be defined by their parameter count alone, but by their ability to dynamically allocate compute—balancing the depth of thought against the constraints of latency, energy, and task complexity.

\section*{References}
\begin{enumerate}\small
  \item Attention Is All You Need (2017). \texttt{arXiv:1706.03762}
  \item Proximal Policy Optimization Algorithms (2017). \texttt{arXiv:1707.06347}
  \item Language Models are Few-Shot Learners (2020). \texttt{arXiv:2005.14165}
  \item Scaling Laws for Neural Language Models (2020). \texttt{arXiv:2001.08361}
  \item Measuring Massive Multitask Language Understanding (2020). \texttt{arXiv:2009.03300}
  \item Learning to summarize from human feedback (2020). \texttt{arXiv:2009.01325}
  \item Evaluating Large Language Models Trained on Code (2021). \texttt{arXiv:2107.03374}
  \item Training Verifiers to Solve Math Word Problems (2021). \texttt{arXiv:2110.14168}
  \item Measuring Mathematical Problem Solving With the MATH Dataset (2021). \texttt{arXiv:2103.03874}
  \item Program Synthesis with Large Language Models (2021). \texttt{arXiv:2108.07732}
  \item Show Your Work: Scratchpads for Intermediate Computation with Language
  Models (2021). \texttt{arXiv:2112.00114}
  \item Chain-of-Thought Prompting Elicits Reasoning in Large Language Models (2022). \texttt{arXiv:2201.11903}
  \item Self-Consistency Improves Chain of Thought Reasoning in Language Models (2022). \texttt{arXiv:2203.11171}
  \item Training language models to follow instructions with human feedback (2022). \texttt{arXiv:2203.02155}
  \item Large Language Models are Zero-Shot Reasoners (2022). \texttt{arXiv:2205.11916}
  \item STaR: Bootstrapping Reasoning With Reasoning (2022). \texttt{arXiv:2203.14465}
  \item Solving Quantitative Reasoning Problems with Language Models (2022). \texttt{arXiv:2206.14858}
  \item Solving math word problems with process- and outcome-based feedback (2022). \texttt{arXiv:2211.14275}
  \item Least-to-Most Prompting Enables Complex Reasoning in Large Language
  Models (2022). \texttt{arXiv:2205.10625}
  \item Training Compute-Optimal Large Language Models (2022). \texttt{arXiv:2203.15556}
  \item Making Large Language Models Better Reasoners with Step-Aware Verifier (2022). \texttt{arXiv:2206.02336}
  \item Training a Helpful and Harmless Assistant with Reinforcement Learning
  from Human Feedback (2022). \texttt{arXiv:2204.05862}
  \item ReAct: Synergizing Reasoning and Acting in Language Models (2022). \texttt{arXiv:2210.03629}
  \item PAL: Program-aided Language Models (2022). \texttt{arXiv:2211.10435}
  \item GPT-4 Technical Report (2023). \texttt{arXiv:2303.08774}
  \item Tree of Thoughts: Deliberate Problem Solving with Large Language Models (2023). \texttt{arXiv:2305.10601}
  \item LLaMA: Open and Efficient Foundation Language Models (2023). \texttt{arXiv:2302.13971}
  \item Self-Refine: Iterative Refinement with Self-Feedback (2023). \texttt{arXiv:2303.17651}
  \item Llama 2: Open Foundation and Fine-Tuned Chat Models (2023). \texttt{arXiv:2307.09288}
  \item Let's Verify Step by Step (2023). \texttt{arXiv:2305.20050}
  \item Math-Shepherd: Verify and Reinforce LLMs Step-by-step without Human
  Annotations (2023). \texttt{arXiv:2312.08935}
  \item GPQA: A Graduate-Level Google-Proof Q\&A Benchmark (2023). \texttt{arXiv:2311.12022}
  \item Efficient Memory Management for Large Language Model Serving with
  PagedAttention (2023). \texttt{arXiv:2309.06180}
  \item Reflexion: Language Agents with Verbal Reinforcement Learning (2023). \texttt{arXiv:2303.11366}
  \item Alphazero-like Tree-Search can Guide Large Language Model Decoding and
  Training (2023). \texttt{arXiv:2309.17179}
  \item Reasoning with Language Model is Planning with World Model (2023). \texttt{arXiv:2305.14992}
  \item Direct Preference Optimization: Your Language Model is Secretly a Reward
  Model (2023). \texttt{arXiv:2305.18290}
  \item MetaMath: Bootstrap Your Own Mathematical Questions for Large Language
  Models (2023). \texttt{arXiv:2309.12284}
  \item Scaling LLM Test-Time Compute Optimally can be More Effective than
  Scaling Model Parameters (2024). \texttt{arXiv:2408.03314}
  \item DeepSeekMath: Pushing the Limits of Mathematical Reasoning in Open
  Language Models (2024). \texttt{arXiv:2402.03300}
  \item The Llama 3 Herd of Models (2024). \texttt{arXiv:2407.21783}
  \item Large Language Monkeys: Scaling Inference Compute with Repeated Sampling (2024). \texttt{arXiv:2407.21787}
  \item OpenAI o1 System Card (2024). \texttt{arXiv:2412.16720}
  \item Qwen2.5 Technical Report (2024). \texttt{arXiv:2412.15115}
  \item Inference Scaling Laws: An Empirical Analysis of Compute-Optimal
  Inference for Problem-Solving with Language Models (2024). \texttt{arXiv:2408.00724}
  \item Do NOT Think That Much for 2+3=? On the Overthinking of o1-Like LLMs (2024). \texttt{arXiv:2412.21187}
  \item The Art of Scaling Test-Time Compute for Large Language Models (2025). \texttt{arXiv:2512.02008}
  \item Scaling over Scaling: Exploring Test-Time Scaling Plateau in Large Reasoning Models (2025). \texttt{arXiv:2505.20522}
  \item Scaling Test-time Compute for LLM Agents (2025). \texttt{arXiv:2506.12928}
  \item The Energy Cost of Reasoning: Analyzing Energy Usage in LLMs with Test-time Compute (2025). \texttt{arXiv:2505.14733}
  \item Reasoning on a Budget: A Survey of Adaptive and Controllable Test-Time Compute in LLMs (2025). \texttt{arXiv:2507.02076}
  \item Inference-Time Scaling for Complex Tasks: Where We Stand and What Lies
  Ahead (2025). \texttt{arXiv:2504.00294}
  \item A Survey on Test-Time Scaling in Large Language Models: What, How,
  Where, and How Well? (2025). \texttt{arXiv:2503.24235}
  \item m1: Unleash the Potential of Test-Time Scaling for Medical Reasoning with Large Language Models (2025). \texttt{arXiv:2504.00869}
  \item Latency and Token-Aware Test-Time Compute (2025). \texttt{arXiv:2509.09864}
  \item Scaling up Test-Time Compute with Latent Reasoning: A Recurrent Depth
  Approach (2025). \texttt{arXiv:2502.05171}
  \item Rethinking Fine-Tuning when Scaling Test-Time Compute: Limiting Confidence Improves Mathematical Reasoning (2025). \texttt{arXiv:2502.07154}
  \item Scaling Evaluation-time Compute with Reasoning Models as Evaluators (2025). \texttt{arXiv:2503.19877}
  \item Are We Scaling the Right Thing? A System Perspective on Test-Time Scaling (2025). \texttt{arXiv:2509.19645}
  \item Inference-Time Computations for LLM Reasoning and Planning: A Benchmark
  and Insights (2025). \texttt{arXiv:2502.12521}
  \item Rethinking Test-Time Scaling for Medical AI: Model and Task-Aware Strategies for LLMs and VLMs (2025). \texttt{arXiv:2506.13102}
  \item Test-Time Scaling in Reasoning Models Is Not Effective for Knowledge-Intensive Tasks Yet (2025). \texttt{arXiv:2509.06861}
  \item On the Role of Temperature Sampling in Test-Time Scaling (2025). \texttt{arXiv:2510.02611}
  \item Prompting Test-Time Scaling Is A Strong LLM Reasoning Data Augmentation (2025). \texttt{arXiv:2510.09599}
  \item Test-time Scaling of LLMs: A Survey from A Subproblem Structure Perspective (2025). \texttt{arXiv:2511.14772}
  \item Enhancing Test-Time Scaling of Large Language Models with Hierarchical Retrieval-Augmented MCTS (2025). \texttt{arXiv:2507.05557}
  \item Limits and Gains of Test-Time Scaling in Vision-Language Reasoning (2025). \texttt{arXiv:2512.11109}
  \item Crosslingual Reasoning through Test-Time Scaling (2025). \texttt{arXiv:2505.05408}
  \item T1: Advancing Language Model Reasoning through Reinforcement Learning and Inference Scaling (2025). \texttt{arXiv:2501.11651}
  \item Adaptive Test-Time Reasoning via Reward-Guided Dual-Phase Search (2025). \texttt{arXiv:2509.25420}
  \item Thought calibration: Efficient and confident test-time scaling (2025). \texttt{arXiv:2505.18404}
  \item ScaleRTL: Scaling LLMs with Reasoning Data and Test-Time Compute for Accurate RTL Code Generation (2025). \texttt{arXiv:2506.05566}
  \item Fractional Reasoning via Latent Steering Vectors Improves Inference Time Compute (2025). \texttt{arXiv:2506.15882}
  \item Can 1B LLM Surpass 405B LLM? Rethinking Compute-Optimal Test-Time
  Scaling (2025). \texttt{arXiv:2502.06703}
  \item AgentTTS: Large Language Model Agent for Test-time Compute-optimal Scaling Strategy in Complex Tasks (2025). \texttt{arXiv:2508.00890}
  \item DeepSeek-R1: Incentivizing Reasoning Capability in LLMs via Reinforcement Learning (2025). \texttt{arXiv:2501.12948}
  \item s1: Simple test-time scaling (2025). \texttt{arXiv:2501.19393}
  \item TMAS: Scaling Test-Time Compute via Multi-Agent Synergy (2026). \texttt{arXiv:2605.10344}
  \item Improving reasoning at inference time via uncertainty minimisation (2026). \texttt{arXiv:2603.07159}
  \item Predicting and improving test-time scaling laws via reward tail-guided search (2026). \texttt{arXiv:2602.01485}
\end{enumerate}

\end{document}
